% D19
% Auto-generated from markdown source
% origin: ai

\documentclass[12pt, a4paper]{article}
\usepackage[utf8]{inputenc}
\usepackage[T1]{fontenc}
\usepackage{amsmath, amssymb, amsfonts, amsthm}
\usepackage{textcomp}
\usepackage{geometry}
\usepackage{hyperref}
\usepackage{booktabs}
\usepackage{graphicx}
\usepackage{fancyhdr}
\usepackage{tcolorbox}
\usepackage{enumitem}

\geometry{margin=1in}
\pagestyle{fancy}
\fancyhf{}
\rhead{D19}
\lhead{\thepage}

\newtheorem{definition}{Definition}[section]
\newtheorem{theorem}{Theorem}[section]
\newtheorem{corollary}{Corollary}[section]
\newtheorem{proposition}{Proposition}[section]
\newtheorem{lemma}{Lemma}[section]
\newtheorem{conjecture}{Conjecture}[section]

\newcommand{\calI}{\mathcal{I}}
\newcommand{\calP}{\mathcal{P}}
\newcommand{\calE}{\mathcal{E}}
\newcommand{\calS}{\mathcal{S}}
\newcommand{\calH}{\mathcal{H}}
\newcommand{\calO}{\mathcal{O}}
\newcommand{\calF}{\mathcal{F}}
\newcommand{\calU}{\mathcal{U}}
\newcommand{\calV}{\mathcal{V}}
\newcommand{\calA}{\mathcal{A}}
\newcommand{\calG}{\mathcal{G}}
\newcommand{\calB}{\mathcal{B}}
\newcommand{\calC}{\mathcal{C}}
\newcommand{\calL}{\mathcal{L}}
\newcommand{\calM}{\mathcal{M}}
\newcommand{\calX}{\mathcal{X}}
\newcommand{\calY}{\mathcal{Y}}
\newcommand{\calW}{\mathcal{W}}
\newcommand{\calT}{\mathcal{T}}
\newcommand{\calZ}{\mathcal{Z}}
\newcommand{\calR}{\mathcal{R}}
\newcommand{\calQ}{\mathcal{Q}}
\newcommand{\calD}{\mathcal{D}}
\newcommand{\calN}{\mathcal{N}}
\newcommand{\calK}{\mathcal{K}}

\newcommand{\Tr}{\text{Tr}}
\newcommand{\E}{\mathbb{E}}
\newcommand{\diff}{\mathrm{d}}
\newcommand{\KL}{D_{\mathrm{KL}}}
\newcommand{\argmax}{\arg\max}
\newcommand{\argmin}{\arg\min}
\newcommand{\hbarF}{\hbar}

\newtcolorbox{thmbox}[1]{
  colback=blue!5!white, colframe=blue!75!black, fonttitle=\bfseries,
  title={#1}, breakable
}
\newtcolorbox{defbox}[1]{
  colback=green!5!white, colframe=green!75!black, fonttitle=\bfseries,
  title={#1}, breakable
}

\begin{document}

\title{D19}
\author{RALPH Loop}
\date{}
\maketitle

\begin{abstract}
This paper presents D19, formalizing the underlying principles using information-theoretic methods.

\textbf{Keywords}: D19, information theory, formalization.
\end{abstract}

\tableofcontents
\section{D19}




\subsection{0.}


****
\begin{itemize}
\item ****
\item ****
\item \textbf{**/}**
\end{itemize}


\textbf{D19 }\textbf{-}****



\subsection{1.}


\subsubsection{1.1}


\textbf{ D19-1 \(\mathcal{L}_{\epsilon}\)}
/**** \(\epsilon\)
\begin{itemize}
\item \textbf{\(\epsilon_{\text{}} \ll 1\)}
\item \textbf{\(\epsilon_{\text{}} \approx \epsilon_0\)}
\item \textbf{\(\epsilon_{\text{}} \gg \epsilon_0\)}
\end{itemize}


****
\[d_{\mathcal{L}}(l_1, l_2) = \frac{1}{\epsilon(l_1)} + \frac{1}{\epsilon(l_2)} - 2 \cdot \min\left(\frac{1}{\epsilon(l_1)}, \frac{1}{\epsilon(l_2)}\right)\]


\subsubsection{1.2}


\textbf{ D19-2 \(\mathcal{C}_{l_1 \to l_2}\)}
\(l_1\)  \(l_2\) ****
\[\mathcal{C}_{l_1 \to l_2}: \mathcal{L}_{\text{}} \times \mathcal{K}_{\text{}} \to \mathcal{L}_{\text{}} \times \mathcal{K}_{\text{}}\]

****
\[C(\mathcal{C}_{l_1 \to l_2}) = \max_{P_{\text{}}} I(\mathcal{K}_{\text{}}; \mathcal{K}'_{\text{}})\]

\(\mathcal{K}'_{\text{}}\)

****
\[\mathcal{L}_{\text{precision}}(l_1, l_2) = \frac{H(\mathcal{K}_{\text{}} | \mathcal{K}'_{\text{}})}{H(\mathcal{K}_{\text{}})}\]


\subsection{2.}


\subsubsection{ D19-1}


****
\((l_H, l_M, l_L)\)

\textbf{D19-1 }
\[C(l_H \to l_L) \leq C(l_H \to l_M) + C(l_M \to l_L) - \Delta_{\text{bridge}}\]

\(\Delta_{\text{bridge}} \geq 0\) ****

\textbf{ D19-1a}
\begin{itemize}
\item ****\(C(l_H \to l_M) + C(l_M \to l_L) \geq C(l_H \to l_L)\)
\item ****\(\Delta_{\text{bridge}} > 0 \iff\)
\end{itemize}



\subsubsection{ D19-2-}


****
\textbf{** \(\eta\) }** \(\alpha\)
\[\eta = \frac{R}{C_{\text{max}}}, \quad \alpha = 1 - \frac{H(\mathcal{K}| \mathcal{K}')}{H(\mathcal{K})}\]

\textbf{D19-2 }
\[\eta \cdot \alpha \leq \frac{1}{2} \log_2\left(1 + \frac{\epsilon_{\text{channel}}}{\sigma^2_{\text{noise}}}\right)\]

\textbf{- Pareto }
\[\mathcal{P}^* = \left\{(\eta, \alpha) \mid \exists p: \eta \cdot \alpha = \max\right\}\]

****
\begin{itemize}
\item \textbf{** }**
\item \textbf{** }**
\item ****\(\frac{\partial(\eta \cdot \alpha)}{\partial p} = 0\)
\end{itemize}



\subsubsection{ D19-3}


****
\textbf{** \(f: \mathcal{K}_{\text{}} \to \mathcal{K}_{\text{}}\)}**
\[\mathcal{Q}(f) = \frac{I(f(\mathcal{K}_{\text{}}); \mathcal{K}_{\text{}})}{H(\mathcal{K}_{\text{}})}\]

\textbf{D19-3 }
\[f^* = \arg\max_{f} \left[\mathcal{Q}(f) - \lambda \cdot \mathcal{D}_{\text{complexity}}(f)\right]\]


\begin{itemize}
\item \(\mathcal{Q}(f)\)
\item \(\mathcal{D}_{\text{complexity}}(f)\)
\item \(\lambda\)
\end{itemize}


****
\[\frac{\delta \mathcal{Q}}{\delta f} = \lambda \cdot \frac{\delta \mathcal{D}_{\text{complexity}}}{\delta f}\]

\textbf{*}***


\subsubsection{ D19-4}


****
**** \(\epsilon_k\) \(\epsilon_{\text{in}}\)

\textbf{D19-4 }
\[\epsilon_{k+1} = \epsilon_k + \Delta \epsilon_k + \eta_k \cdot \epsilon_{\text{in}}\]


\begin{itemize}
\item \(\Delta \epsilon_k\) \(k\) ****
\item \(\eta_k\) \(k\) ****
\end{itemize}


****
\[\epsilon_N = \epsilon_0 \cdot \prod_{k=0}^{N-1} (1 + \eta_k) + \sum_{k=0}^{N-1} \Delta \epsilon_k \cdot \prod_{j=k+1}^{N-1} (1 + \eta_j)\]

\textbf{ D19-4a}
\begin{itemize}
\item ****\(\Delta \epsilon_k \approx \text{machine epsilon}\)
\item \textbf{/ML }\(\Delta \epsilon_k\)
\end{itemize}


**** =


\subsubsection{ D19-5}


****
\textbf{*}***
\[\mathcal{I}_{\text{conserved}}(\mathcal{K}_{l_1 \to l_2}) = \frac{I(\mathcal{K}_{l_1}; f(\mathcal{K}_{l_2})))}{H(\mathcal{K}_{l_1})}\]

\textbf{D19-5 }
\[\mathcal{I}_{\text{conserved}}(l_H \to l_L) \geq \mathcal{I}_{\text{conserved}}(l_H \to l_M) \cdot \mathcal{I}_{\text{conserved}}(l_M \to l_L)\]

\textbf{ D19-5a}
\[\mathcal{I}_{\text{conserved}}(l_H \to l_L) \leq \mathcal{I}_{\text{conserved}}(l_H \to l_M) \cdot \mathcal{I}_{\text{conserved}}(l_M \to l_L)\]

\textbf{*}\textbf{*   \(f_{HM}\)  \(f_{ML}\) }**


\subsection{3.}


\subsubsection{ D19-A}


\textbf{ D19-5}
1.
2.  vs  vs
3.
4. **** >


\subsubsection{ D19-B- Pareto}


\textbf{ D19-2}
1.
2.
3.  Pareto
4. ****


\subsubsection{ D19-C}


\textbf{ D19-4}
1. FP32FP16INT8
2.  \(\epsilon_k\)
3. ****


\subsubsection{ D19-D}


\textbf{ D19-3}
1.
2.  \(\lambda\)
3. -
4. **** \(\lambda^*\)


\subsection{4.}


\subsubsection{4.1 D16}


\textbf{D19-D16 }
\begin{itemize}
\item D19  \(\epsilon\)  D16 ****
\item D19-2 - = D16-4 ****
\item \textbf{** = }**
\end{itemize}



\subsubsection{4.2 D7}


\textbf{D19-D7 }
\begin{itemize}
\item D19  \(\Delta \epsilon_k\) = D7 ****
\item = ****
\item \textbf{** = }**
\end{itemize}



\subsubsection{4.3 D13}


\textbf{D19-D13 }
\begin{itemize}
\item D19  \(\mathcal{C}_{l_1 \to l_2}\)  D13 ****
\item = D13 ****
\item \textbf{** = }**
\end{itemize}



\subsection{5.}


\subsubsection{5.1}


1. \textbf{D19-1 }
2. \textbf{D19-2 }
3. \textbf{D19-4 }
4. \textbf{D19-5 }

\subsubsection{5.2}


1. \textbf{/}
2. ****
3. ****


\subsection{6.}


###

**** \(\mathcal{L}_{\text{}}\) \(\mathcal{L}_{\text{}}\) \(\mathcal{L}_{\text{}}\) \(\epsilon_H \ll \epsilon_M \ll \epsilon_L\)

****
1. ****\(C(l_H \to l_L) \leq C(l_H \to l_M) + C(l_M \to l_L)\)
2. ****\(f^* = \arg\max_f [\mathcal{Q}(f) - \lambda \cdot \mathcal{D}(f)]\)
3. ****\(\epsilon_N = \epsilon_0 \prod(1 + \eta_k) + \sum \Delta \epsilon_k \prod(1 + \eta_j)\)
4. ****\(\mathcal{I}_{\text{conserved}}(l_H \to l_L) = \mathcal{I}_{\text{conserved}}(l_H \to l_M) \cdot \mathcal{I}_{\text{conserved}}(l_M \to l_L)\)

\textbf{**  }\textbf{**Pareto }


*D19 -*

\begin{thebibliography}{99}
\bibitem{ref1} Shannon, C.E. (1948). A Mathematical Theory of Communication.
\bibitem{ref2} Cover, T.M. (2006). Elements of Information Theory.
\end{thebibliography}

\end{document}